\section{The coefficient test and soundness}
\label{sec:coefficient}

The remaining task is to bound the density of an independent set after
fixing an aligned pair of a chain and a selected set of layers.  We add
one auxiliary phase for each selected layer.  Alignment makes the
influence of a zero-coefficient layer small, whereas the geometric
coefficient distribution makes large influences at positive coefficients
unlikely.

Keep the tolerance $\eps>0$ of Section~\ref{sec:alignment}, and take
$\beta>0$ and $K\ge1$ from
Lemma~\ref{lem:influences} with $L=64$ and this $\eps$.
For the graph parameters of Section~\ref{sec:construction} and the
list-approximation tolerance $\gamma$ of Section~\ref{sec:lists}, assume
\begin{equation}\label{eq:parameter-bounds}
 \begin{gathered}
 L/m<\beta/2,\qquad \lambda K<\eps,
 \qquad (1-p_m)^s<\eps,\\
 \frac{s(2\gamma)^2}{\beta^2}<\eps,\qquad
 \frac{2L}{D}<\gamma,\qquad \frac{2Ls}{P}<\eps.
 \end{gathered}
\end{equation}
Section~\ref{sec:parameters} will choose constants satisfying these
inequalities. A coefficient equal to $1$ is called a \emph{full
coefficient}. Although its individual probability $p_m$ may be small,
the third inequality makes at least one of the $s$ coefficients full
with probability greater than $1-\eps$.

\begin{proposition}\label{prop:aligned-bound}
Assume \eqref{eq:parameter-bounds}. Suppose that the construction does not
take the clique branch, and let
$\mathcal A$ be a nonempty independent set of its output graph.  Form the
functions and lists associated with $\mathcal A$ as above.  For every
chain and $s$-element set $B$ admitting a map
\[
 a:M_r\longrightarrow B,
 \qquad a(A_I)\subseteq I
 \quad\text{for every nonempty }I\subseteq B,
\]
the conditional probability that the sampled output vertex belongs to
$\mathcal A$ is at most $6\eps$.
\end{proposition}

\begin{proof}
Fix such a chain, $B$, and a map $a$, before drawing any coefficients or
rotations.  Write $b=\min B$.
Since $f=1$ at every location of $\mathcal A$, it suffices to bound the
expected square of $F_b$ at the sampled grid location by $6\eps$.
We establish this through a continuous rotation experiment with auxiliary
phases.  Initially replace the discrete rotations
by independent uniform real variables
\[
 \theta_{j,\ell}\in[0,1),
 \qquad j\in B,\quad \ell\in M_j.
\]
All these variables are independent of the coefficient vector
$t\in\mathcal T_m^B$.  Introduce additional independent Haar-uniform
variables $z=(z_i)_{i\in B}\in\T^B$, and define
\begin{equation}\label{coeff:H}
 H(t,\theta;z)
 =F_b\left(\left(
 z_{a(\pi_{br}(\kappa))}
 +\sum_{j\in B}t_j\theta_{j,\pi_{bj}(\kappa)}
 \pmod 1\right)_{\kappa\in M_b}\right).
\end{equation}
For fixed $t,\theta$, the map from $z$ to the phase vector in
\eqref{coeff:H} is nonexpanding for the sup circle metrics: each output
coordinate uses just one coordinate of $z$.  Consequently $H$ is
$L$-Lipschitz as a function of $z$, independently of $s$ and of all label
cardinalities.  It is bounded by one in absolute value.  A simultaneous
half-shift of the $z$ coordinates half-shifts every input coordinate of
$F_b$, so oddness gives
\begin{equation}\label{coeff:mean-zero}
 \E_z H(t,\theta;z)=0.
\end{equation}

Let $E_i$ average only $z_i$ and let $Q_i=\id-E_i$, an orthogonal
projection of norm one.  We use the notation
\[
 D_i(t,\theta)=\|Q_iH(t,\theta;\cdot)\|_{L^2(\T^B)},
 \qquad X(t,\theta)=\E_z H(t,\theta;z)^2.
\]
By Lemma~\ref{lem:influences} and \eqref{coeff:mean-zero},
\begin{align}
 X(t,\theta)>\eps
 &\quad\Longrightarrow\quad
 \max_{i\in B}D_i(t,\theta)\ge\beta,
 \label{coeff:variance-detection}\\
 \#\{i\in B:D_i(t,\theta)\ge\beta/2\}
 &\le K
 \quad\text{for every }t,\theta.
 \label{coeff:count}
\end{align}

\paragraph{Zero coefficients.}
Let
\[
 \mathcal F=\{t:\text{some }t_j=1\}
\]
be the event that a full coefficient is present.  Fix $t\in\mathcal F$
and an index $i\in B$ with $t_i=0$.  Choose $j\in B$ with $t_j=1$;
necessarily $j\ne i$.  Put $I=B\setminus\{i\}$ and $c=\min I$.
The set $I$ is nonempty.  In the subchain function
$h_{I,t|_I}$ from \eqref{eq:subchain-functions}, leave all rotation
blocks unchanged except block $j$, where we substitute
\begin{equation}\label{coeff:shift}
 \theta'_{j,\ell}
 =\theta_{j,\ell}+z_{a(\pi_{jr}(\ell))}\pmod 1,
 \qquad \ell\in M_j.
\end{equation}
Here and below a phase used as a real rotation is represented in
$[0,1)$.  For each fixed $z$, this substitution preserves the product
uniform distribution of the rotation inputs.  Even when several
coordinates receive the same shift, each coordinate undergoes a fixed
translation, so their conditional joint law remains the same product
law.

To compare the two function evaluations, let $y\in\T^{M_c}$ be the
phase vector in $h_{I,t|_I}(\theta')$.  For every $\kappa\in M_b$,
composition of the projections gives
\begin{equation}\label{coeff:pullback-identity}
 y_{\pi_{bc}(\kappa)}
 =z_{a(\pi_{br}(\kappa))}
   +\sum_{k\in B}t_k\theta_{k,\pi_{bk}(\kappa)}
   \pmod 1.
\end{equation}
Indeed, the omitted $i$ term is zero, and the changed $j$ term has
coefficient exactly one.  This last fact is essential: reducing the
shifted rotation modulo one before multiplying would not in general be
valid at a fractional coefficient.  Thus
$H(t,\theta;z)=F_b(y\circ\pi_{bc})$.  If $c=b$ the two evaluations
are identical; if deleting $i$ changes the minimum layer, the
compatibility bound \eqref{eq:compatibility} still gives
\begin{equation}\label{coeff:deletion-error}
 |H(t,\theta;z)-h_{I,t|_I}(\theta')|
 \le 2L/D<\gamma.
\end{equation}

Let $g_{I,t|_I}$ be the fixed junta approximating $h_{I,t|_I}$ in
$L^2$ norm to error less than $\gamma$.  Its selected coordinates in
block $k$ belong to $S_I^k$.  Product-law preservation for every fixed
$z$ implies
\[
 \big\|h_{I,t|_I}(\theta')-g_{I,t|_I}(\theta')\big\|_{L^2(\theta,z)}
 <\gamma.
\]
Moreover, the composed function $g_{I,t|_I}(\theta')$ is pointwise
independent of $z_i$.  Only the modified block $j$ introduces any $z$
dependence.  Every selected coordinate $\ell\in S_I^j$ satisfies
$\pi_{jr}(\ell)\in A_I$ by \eqref{eq:terminal-lists}, and alignment
therefore gives $a(\pi_{jr}(\ell))\in I$.  Combining this observation
with \eqref{coeff:deletion-error} and applying the contraction $Q_i$
on the full product space yields
\begin{equation}\label{coeff:zero-bound}
 \E_\theta D_i(t,\theta)^2
 =\|Q_iH\|_{L^2(\theta,z)}^2
 \le
 \big\|H-g_{I,t|_I}(\theta')\big\|_{L^2(\theta,z)}^2
 <(2\gamma)^2.
\end{equation}
This estimate holds for every fixed $t\in\mathcal F$ with $t_i=0$.
Markov's inequality, followed by a union bound over the $s$ indices,
therefore gives
\begin{equation}\label{coeff:zero-probability}
 \Pr_{t,\theta}\bigl(
 t\in\mathcal F,\ 
 \exists i\in B:\ t_i=0,\ D_i(t,\theta)\ge\beta
 \bigr)
 \le \frac{s(2\gamma)^2}{\beta^2}<\eps.
\end{equation}

\paragraph{Positive coefficients.}
We next bound large influences at positive coefficient levels.  This
step does not assume that a full coefficient is present.  Fix $i$,
all rotations $\theta$, and all coefficients other than $t_i$.
For $0\le k\le m$, let $H_k$ denote \eqref{coeff:H} with $t_i=k/m$,
and put $d_k=\|Q_iH_k\|_2$.  Changing one coefficient by $1/m$
changes each phase coordinate by circle distance at most $1/m$, since
the rotations lie in $[0,1)$.  The norm-one property of $Q_i$ gives
\begin{equation}\label{coeff:one-step}
 |d_k-d_{k-1}|
 \le\|Q_i(H_k-H_{k-1})\|_2
 \le L/m<\beta/2
 \qquad(1\le k\le m).
\end{equation}
Write $p_k=\lambda^k/\sum_{h=0}^m\lambda^h$ for the probability
of coefficient level $k/m$.  Since $p_k=\lambda p_{k-1}$,
\eqref{coeff:one-step} implies, for the fixed background data,
\begin{align*}
 \sum_{k=1}^m p_k\mathbf 1_{\{d_k\ge\beta\}}
 &\le\lambda\sum_{h=0}^{m-1}
              p_h\mathbf 1_{\{d_h\ge\beta/2\}}\\
 &\le\lambda\sum_{h=0}^m
              p_h\mathbf 1_{\{d_h\ge\beta/2\}}.
\end{align*}
Integrating over the fixed data and summing over $i\in B$, we obtain
from \eqref{coeff:count}
\begin{equation}\label{coeff:positive-probability}
 \begin{split}
 \Pr_{t,\theta}\bigl(
 \exists i\in B:\ t_i>0,\ D_i(t,\theta)\ge\beta
 \bigr)
 &\le\lambda\sum_{i\in B}
          \Pr_{t,\theta}(D_i(t,\theta)\ge\beta/2)\\
 &\le\lambda K<\eps.
 \end{split}
\end{equation}
Decreasing a level can destroy the only full coefficient.  This causes
no restriction here: \eqref{coeff:count} holds for every coefficient
vector, including all images of the one-level decrease.

\paragraph{From influences to mean square.}
The parameter choices in \eqref{eq:parameter-bounds} give
$\Pr(t\notin\mathcal F)<\eps$.  By
\eqref{coeff:variance-detection}, the event $X(t,\theta)>\eps$ is
contained in the union of this event and the two events bounded in
\eqref{coeff:zero-probability} and
\eqref{coeff:positive-probability}.  Consequently
\[
 \Pr_{t,\theta}(X(t,\theta)>\eps)<3\eps.
\]
Since $0\le X\le1$, it follows that
\begin{equation}\label{coeff:mean-square}
 \E_{t,\theta,z}H(t,\theta;z)^2
 =\E_{t,\theta}X(t,\theta)
 \le\eps+\Pr(X>\eps)<4\eps.
\end{equation}
We now remove the auxiliary phases and return to the actual grid
sampling.  This is the only remaining passage from the analytic test
to the conditional vertex density.

\paragraph{Removing the phases and discretizing.}
Define the continuous location
\[
 x(t,\theta)
 =\left(\sum_{j\in B}t_j\theta_{j,\pi_{bj}(\kappa)}
          \pmod 1\right)_{\kappa\in M_b}.
\]
For every fixed $t\in\mathcal F$, choose an index $j$ with $t_j=1$.
Applying the translation \eqref{coeff:shift} to this block absorbs
every $z$ term in \eqref{coeff:H}.  For each fixed $z$ the translated
rotation array again has exactly its original product law.  Hence
\begin{equation}\label{coeff:absorption}
 \E_{\theta,z}H(t,\theta;z)^2
 =\E_\theta F_b(x(t,\theta))^2
 \qquad(t\in\mathcal F).
\end{equation}
The choice of the full block may depend on $t$, because this equality
is asserted separately for every fixed coefficient vector.  In
particular, \eqref{coeff:mean-square} implies
\begin{equation}\label{coeff:continuous-bound}
 \E_{t,\theta}\!\left[
 \mathbf 1_{\mathcal F}(t)F_b(x(t,\theta))^2\right]<4\eps.
\end{equation}

Couple the discrete rotations to the continuous ones by
\[
 \widehat\theta_{j,\ell}
 =\frac{\lfloor P\theta_{j,\ell}\rfloor}{P}.
\]
They have the required independent uniform grid law.  Each input
coordinate changes by less than $1/P$, so the sup circle distance
between $x(t,\theta)$ and $x(t,\widehat\theta)$ is at most $s/P$.
The Lipschitz bound and the range $[-1,1]$ therefore imply
\begin{equation}\label{coeff:grid-error}
 \left|F_b(x(t,\theta))^2
       -F_b(x(t,\widehat\theta))^2\right|
 \le 2Ls/P<\eps.
\end{equation}
Combining \eqref{coeff:continuous-bound},
\eqref{coeff:grid-error}, and $\Pr(t\notin\mathcal F)<\eps$ gives
\begin{equation}\label{coeff:discrete-bound}
 \E_{t,\widehat\theta}
       F_b(x(t,\widehat\theta))^2
 <4\eps+\eps+\eps=6\eps.
\end{equation}
The discrete location belongs to the grid of denominator $D=mP$, so
$F_b$ equals the original grid function $f$ there.  That function
equals one at every location of a vertex of $\mathcal A$.
Thus the indicator that the sampled vertex belongs to $\mathcal A$
is pointwise bounded by $F_b(x(t,\widehat\theta))^2$.
This remains true when several distinct vertices have the same
location.  Inequality~\eqref{coeff:discrete-bound} proves the
proposition.
\end{proof}
